\documentclass[a5paper,10pt]{article}

% https://github.com/InDevRus/filippov-solutions

% Нумерация страниц
\pagenumbering{gobble}

% Настройка полей
\usepackage{geometry}
\geometry{left=1cm}
\geometry{right=1cm}
\geometry{top=1cm}
\geometry{bottom=1.5cm}

% Вставка таблицы
\usepackage{array}
\usepackage{tabularx}

% Инструменты выравнивания формул
\usepackage{amsmath}
\usepackage{mathtools}

% Диакритика в математике
\usepackage{amsfonts}

% Вычисления размеров на ходу
\usepackage{calc}

% Удобные записи производных
\usepackage[italic=true]{derivative}

% Настройки сносок
\usepackage{enotez}

% Длинные знаки векторов
\usepackage{anyfontsize}
\usepackage[b]{esvect}

% Вставка картинок
\usepackage{graphicx}

% Вставка красной строки
\usepackage{indentfirst}
\setlength{\parindent}{1em}

% Условия в командах
\usepackage{xifthen}

% Луа-скрипты
\usepackage{luacode}

% Настройка команд отдельно в математическом и текстовом режимах
\usepackage{mathcommand}

% Для повторения LaTeX кода
\usepackage{multido}

% Конфигурация отступов формул
\delimitershortfall-1sp
\usepackage{mleftright}
\mleftright

% Растягивание объектов
\usepackage{relsize}

% Обтекание текстом
\usepackage{wrapfig}
\usepackage{subcaption}
\usepackage{floatflt}

% Поддержка ввода кириллицы
\usepackage[english, russian]{babel}

% Настройка корневой позиции для компилятора
\usepackage{import}
\providecommand*{\ProjectRootPath}{..}

\import{\ProjectRootPath/preambles/macros/}{theorems}

\import{\ProjectRootPath/preambles/fonts}{font_setup}

% Знак градуса
\usepackage{siunitx}

\import{\ProjectRootPath/preambles/macros/}{operators}
\import{\ProjectRootPath/preambles/macros/}{redefinitions}

\import{\ProjectRootPath/preambles/fonts}{letters_setup}
\import{\ProjectRootPath/preambles/fonts/adjustments}{custom_superscripts}

\import{\ProjectRootPath/preambles/custom}{formatting}
\import{\ProjectRootPath/preambles/custom}{frames}
\import{\ProjectRootPath/preambles/custom}{list_setup}

% TODO: перевести команды на современный стиль объявления

\begin{document}

Неограниченное при $t \to +\infty$ решение обозначим $\Psi\br{t}$,
нулевое -- $\Phi\br{t}$. 
Поскольку это решение линейной однородной системы, всякий столбец
$k \Psi\br{t}$, $k \ne 0$ также является решением, при этом неограниченным. 
Найдём точку $\sub{t}{0}$ такую, что $\vbr{\Psi\br{\sub{t}{0}}} \ne 0$.

Положим $\epsilon = 1$, $\delta > 0$ -- произвольное число и 
$\sub{\Psi}{\delta}\br{t} = \dfrac {\delta \Psi\br{t}} {2\vbr{\Psi\br{\sub{t}{0}}}}$
-- частное решение.
Для такого решения выполняется
$\vbr{\sub{\Psi}{\delta}\br{\sub{t}{0}} - \Phi\br{t}}
    = \vbr{\sub{\Psi}{\delta}\br{\sub{t}{0}}} = \dfrac {\delta} {2} < \delta$.
Этот столбец является неограниченным решением, значит с выбранными 
$\epsilon$ и $\delta$ существует такое число $\sub{t}{1} > \sub{t}{0}$,
что 
$\vbr{\sub{\Psi}{\delta}\br{\sub{t}{1}} - \Phi\br{\sub{t}{1}}}
    = \vbr{\sub{\Psi}{\delta}\br{\sub{t}{1}}} \ge 1 = \epsilon$.
Неустойчивость нулевого решения показана.

\end{document}
